\documentclass[11pt,a4paper]{article}

\usepackage[a4paper,margin=18mm,headheight=24pt]{geometry}
\usepackage[T1]{fontenc}
\usepackage{mathpazo}
\usepackage[scaled=0.92]{helvet}
\usepackage[english]{babel}
\usepackage{microtype}
\usepackage{amsmath,mathtools}
\usepackage{booktabs,tabularx,array}
\usepackage{framed}
\usepackage{fancyhdr}
\usepackage{enumitem}
\usepackage{hyperref}
\usepackage{xcolor}

\definecolor{navy}{HTML}{17365D}
\definecolor{blue}{HTML}{2F75B5}
\definecolor{sky}{HTML}{EAF3FA}
\definecolor{gold}{HTML}{D69E2E}
\definecolor{ink}{HTML}{243447}
\definecolor{muted}{HTML}{617285}
\color{ink}

\hypersetup{colorlinks=true,linkcolor=blue,urlcolor=blue,
	pdftitle={Transformer Model Formula Sheet}}
\pagestyle{fancy}
\fancyhf{}
\lhead{\sffamily\color{muted}Transformer model}
\rhead{\sffamily\color{muted}Formula sheet}
\cfoot{\sffamily\color{muted}\thepage}
\renewcommand{\headrulewidth}{0.4pt}
\renewcommand{\headrule}{\hbox to\headwidth{\color{blue}\leaders\hrule height \headrulewidth\hfill}}

% Lightweight boxes using only framed and xcolor.
\newenvironment{formulaBox}[1]{%
	\par\medskip
	{\sffamily\bfseries\large\color{navy}#1}\par\nobreak\smallskip
	\def\FrameCommand{\hspace{1pt}{\color{blue}\vrule width 2.2pt}\hspace{8pt}}%
	\MakeFramed{\advance\hsize-11.2pt\FrameRestore}%
}{\endMakeFramed\par\smallskip}
\definecolor{shadecolor}{HTML}{F1F6FA}
\newenvironment{noteBox}{%
	\par\smallskip
	\def\FrameCommand{\hspace{1pt}{\color{gold}\vrule width 2.2pt}%
		\colorbox{shadecolor}}%
	\MakeFramed{\advance\hsize-2\width\advance\hsize-2\fboxsep\FrameRestore}%
}{\endMakeFramed\par\smallskip}
\newcommand{\diff}{\mathop{}\!\mathrm{d}}
\newcommand{\Rmag}{\mathcal{R}}
\setlength{\parindent}{0pt}
\setlength{\parskip}{5pt}
\setlist[itemize]{leftmargin=5mm,itemsep=1pt,topsep=2pt}

\begin{document}
	
	\begingroup
	\setlength{\fboxsep}{8mm}%
	\noindent\colorbox{navy}{%
		\parbox{\dimexpr\textwidth-2\fboxsep\relax}{%
			{\sffamily\bfseries\fontsize{25}{29}\selectfont\color{white}
				Transformer Model}\par
			\vspace{2mm}
			{\sffamily\fontsize{15}{18}\selectfont\color{white!85}
				Formula sheet for a three-phase two-winding transformer}%
	}}
	\endgroup
	
	\vspace{3mm}
	\begin{noteBox}
		\textbf{Notation.} Indices $k=1,2,3$ identify the transformer legs or phases;
		$p$ and $s$ denote the primary and secondary windings. Uppercase $U$ and $I$
		denote electrical quantities, $\Phi$ the magnetic flux, $\Psi=N\Phi$ the flux
		linkage, and $\Theta=NI$ the magnetomotive force. The example uses vector group
		$\mathrm{Yy0}$.
	\end{noteBox}
	
	\begin{formulaBox}{Transformer-leg equations}
		For each leg $k\in\{1,2,3\}$, the reluctance associated with the iron-loss
		resistance is
		\begin{equation}
			\Rmag_{R\mathrm{Fe}}=\frac{N_1^2}{R_{\mathrm{Fe}}}.
		\end{equation}
		The nonlinear magnetizing characteristic is represented by
		\begin{equation}
			\Theta_{\mu k}=N_1 I_{\mu N}
			\left[
			a\frac{\Phi_k}{\Phi_N}
			+(1-a)\left(\frac{\Phi_k}{\Phi_N}\right)^m
			\right],
			\qquad k=1,2,3.
		\end{equation}
		Equivalently, the magnetizing current of one leg is
		\begin{equation}
			I_{\mu k}=I_{\mu N}
			\left[
			a\frac{\Phi_k}{\Phi_N}
			+(1-a)\left(\frac{\Phi_k}{\Phi_N}\right)^m
			\right].
		\end{equation}
	\end{formulaBox}
	
	\begin{formulaBox}{Primary-winding equations}
		For every phase $k$,
		\begin{align}
			\Phi_{kp}
			&=\frac{\Psi_{kp}}{N_1}
			=\frac{1}{N_1}\int\!\left(U_{kp}-R_1I_{kp}\right)\diff t,
			& I_{kp}&=\frac{\Theta_k}{N_1}.
		\end{align}
		Written separately for the three phases, this relation applies for
		$k=1$, $2$, and $3$ with the corresponding phase quantities.
	\end{formulaBox}
	
	\begin{formulaBox}{Secondary-winding equations}
		For every phase $k$,
		\begin{align}
			\Phi_{ks}
			&=\frac{\Psi_{ks}}{N_2}
			=\frac{1}{N_2}\int\!\left(U_{ks}-R_2I_{ks}\right)\diff t,
			& I_{ks}&=\frac{\Theta_k}{N_2}.
		\end{align}
		The voltage and current reference directions must be chosen consistently with
		the winding polarities used in the circuit model.
	\end{formulaBox}
	
	\newpage
	\thispagestyle{fancy}
	
	\begin{formulaBox}{Common core flux and winding flux linkages}
		For an ideal magnetic coupling of the windings on each leg,
		\begin{equation}
			\Phi_{kp}=\Phi_{ks}=\Phi_k,
			\qquad k=1,2,3.
		\end{equation}
		The corresponding flux linkages are
		\begin{equation}
			\Psi_{kp}=N_1\Phi_k,
			\qquad
			\Psi_{ks}=N_2\Phi_k.
		\end{equation}
		Thus, in general,
		\begin{equation}
			\Psi_{kp}\ne\Psi_{ks}\quad\text{when}\quad N_1\ne N_2,
		\end{equation}
		even though both windings link the same core flux $\Phi_k$.
	\end{formulaBox}
	
	\begin{formulaBox}{Leakage-flux branch}
		With $\Rmag_{\Phi\sigma}$ denoting the equivalent magnetic reluctance of the
		leakage path,
		\begin{equation}
			\Phi_{k\sigma}=\frac{\Theta_k}{\Rmag_{\Phi\sigma}},
			\qquad k=1,2,3.
		\end{equation}
		For a leakage inductance $L_{1\sigma}$ referred to the primary winding,
		\begin{equation}
			L_{1\sigma}=\frac{N_1^2}{\Rmag_{\Phi\sigma}}
			\qquad\Longleftrightarrow\qquad
			\Rmag_{\Phi\sigma}=\frac{N_1^2}{L_{1\sigma}}.
		\end{equation}
	\end{formulaBox}
	
	\begin{formulaBox}{Tank representation}
		The tank is represented by an additional magnetic reluctance $\Rmag_{\mathrm{tank}}$.
		The source text recommends the following approximate choices:
		\begin{align}
			\Rmag_{\mathrm{tank}}
			&\approx (0.1\ldots0.01)\frac{N_1^2}{L_{1\sigma}},
			&&\text{rough initial estimate - (~10x-100x leakage inductance)},\\
			\Rmag_{\mathrm{tank}}
			&\approx 10^5\ldots10^6,
			&&\text{tank magnetic reluctance in the model units},\\
			\Rmag_{\mathrm{tank}}
			&\approx 10^{20},
			&&\text{tank branch absent (practically open)}.
		\end{align}
	\end{formulaBox}
	
	
	\newpage
	\thispagestyle{fancy}
	
	\begin{formulaBox}{Rated data: 300 MVA, 400 kV / 120 kV, 50 Hz}
		\renewcommand{\arraystretch}{1.18}
		\begin{tabularx}{\textwidth}{>{\bfseries\(}l<{\)} >{\(}l<{\)} X}
			\toprule
			\text{Symbol} & \text{Value} & \textbf{Meaning} \\
			\midrule
			S_N & 300\times10^6\ \mathrm{VA} & Three-phase rated apparent power \\
			U_{1N,LL} & 400\times10^3\ \mathrm{V} & Primary rated line-to-line voltage \\
			U_{2N,LL} & 120\times10^3\ \mathrm{V} & Secondary rated line-to-line voltage \\
			f_N & 50\ \mathrm{Hz} & Rated frequency \\
			N_1 & 1000 & Primary number of turns \\
			u_R & 0.002 & Resistive part of relative short-circuit voltage \\
			u_X & 0.2 & Inductive part of relative short-circuit voltage \\
			i_0 & 0.0006 & Relative no-load current \\
			p_0 & 0.0004 & Relative no-load loss, with $p_0<i_0$ \\
			a & 0.75 & Linear weighting in the magnetizing characteristic \\
			m & 9 & Exponent of the nonlinear magnetizing term \\
			\bottomrule
		\end{tabularx}
	\end{formulaBox}
	
	\begin{formulaBox}{Per-leg quantities for vector group $\mathrm{Yy0}$}
		The three-phase values are first converted to one transformer leg:
		\begin{align}
			S_{N,\mathrm{leg}} &= \frac{S_N}{3}, &
			U_{1N} &= \frac{U_{1N,LL}}{\sqrt{3}}, &
			U_{2N} &= \frac{U_{2N,LL}}{\sqrt{3}},\\
			I_{1N} &= \frac{S_{N,\mathrm{leg}}}{U_{1N}}, &
			Z_{1N} &= \frac{U_{1N}}{I_{1N}}, &
			\omega_N &= 2\pi f_N.
		\end{align}
	\end{formulaBox}
	
	\begin{formulaBox}{Winding and iron-loss parameters}
		The short-circuit impedance is split equally between the two windings and
		referred initially to the primary side:
		\begin{align}
			R_1 &= \frac{Z_{1N}u_R}{2}, &
			L_{1\sigma} &= \frac{Z_{1N}u_X}{2\omega_N},\\
			R_{\mathrm{Fe}} &= \frac{U_{1N}^2}{p_0S_{N,\mathrm{leg}}}, &
			\Rmag_{R\mathrm{Fe}} &= \frac{N_1^2}{R_{\mathrm{Fe}}},\\
			N_2 &= N_1\frac{U_{2N}}{U_{1N}}, &
			R_2 &= R_1\left(\frac{N_2}{N_1}\right)^2,\\
			\Rmag_{\Phi\sigma} &= \frac{N_1^2}{L_{1\sigma}}. &&
		\end{align}
	\end{formulaBox}
	
	\newpage
	\thispagestyle{fancy}
	
	\begin{formulaBox}{Magnetizing branch and rated flux}
		From the no-load current and no-load losses,
		\begin{align}
			X_h &= \frac{U_{1N}^2}
			{S_{N,\mathrm{leg}}\sqrt{i_0^2-p_0^2}},\\
			I_\mu &= \frac{U_{1N}}{X_h},\\
			I_{\mu N} &= \sqrt{2}\,I_\mu,\\
			\Phi_N &= \frac{\sqrt{2}\,U_{1N}}{\omega_NN_1}.
		\end{align}
		The nonlinear characteristic then becomes
		\begin{equation}
			I_\mu(\Phi)=I_{\mu N}
			\left[
			a\frac{\Phi}{\Phi_N}
			+(1-a)\left(\frac{\Phi}{\Phi_N}\right)^m
			\right].
		\end{equation}
	\end{formulaBox}
	
	\begin{formulaBox}{Parameter definitions}
		\renewcommand{\arraystretch}{1.18}
		\begin{tabularx}{\textwidth}{>{\bfseries\(}l<{\)} X}
			\toprule
			\text{Symbol} & \textbf{Meaning} \\
			\midrule
			\Theta_k & Magnetomotive force of leg $k$ \\
			\Phi_k,\ \Phi_{k\sigma} & Main and leakage fluxes of leg $k$ \\
			\Psi_{kp},\ \Psi_{ks} & Primary and secondary flux linkages \\
			\Rmag_{R\mathrm{Fe}} & Magnetic reluctance associated with the iron-loss branch \\
			\Rmag_{\Phi\sigma} & Equivalent leakage-path reluctance \\
			\Rmag_{\mathrm{tank}} & Tank-path magnetic reluctance \\
			R_1,\ R_2 & Primary and secondary winding resistances \\
			R_{\mathrm{Fe}} & Equivalent iron-loss resistance \\
			L_{1\sigma} & Primary-referred leakage inductance \\
			X_h & Main magnetizing reactance \\
			I_{\mu N},\ \Phi_N & Rated peak magnetizing current and rated peak flux \\
			a,\ m & Shape parameters of the nonlinear magnetization curve \\
			\bottomrule
		\end{tabularx}
	\end{formulaBox}
	
	\begin{noteBox}
		\textbf{Implementation notes.} The time integrals may be implemented with
		behavioural sources in LTspice. When forming the winding equations, verify the
		chosen voltage polarities, current reference directions, initial flux, and the
		sign of each magnetomotive-force contribution. The numerical tank-reluctance
		range is model- and unit-dependent and should be validated against the intended
		zero-sequence behaviour.
	\end{noteBox}
	
	\vfill
	{\small\color{muted}\sffamily
		Typeset from the supplied transformer-model notes. English wording and notation
		have been standardized; explicitly identified dimensional inconsistencies in
		the source have been corrected.}
	
\end{document}
